\bookchapter{第二十六章：當下手機技術前沿}

\chapterintro{聊完了舊手機的“身後事”，讓我們把目光拉回當下。2024-2026年是手機技術的一個集中爆發期：AI從雲端塞進了口袋，衛星通訊從旗艦下放到了中端機，摺疊屏從“嚐鮮”變成了“常用”，電池技術也迎來了多年未見的材料級突破。這一章，我們就來盤點這些正在發生的技術前沿，幫你判斷“現在該不該換機”以及“換什麼機”。}

\booksection{一、AI手機：是營銷噱頭，還是真變革？}

2024年起，幾乎所有安卓廠商都在旗艦機上打出“AI手機”的旗號，蘋果也在iOS 18中上線了Apple Intelligence。但到底什麼樣的手機才算“AI手機”？業內大致形成了三條判斷標準：

\begin{itemize}
  \item \textbf{端側大模型}：手機本地能跑引數量在7B左右的語言模型（如Qwen、Llama、DeepSeek的輕量版本），不依賴雲端就能完成文字摘要、翻譯、文案生成等任務。判斷方法：飛航模式下若還能用AI功能，則說明是端側模型。
  \item \textbf{NPU算力（TOPS）}：NPU（神經網路處理單元）的算力以TOPS（每秒萬億次操作）為單位。2024年旗艦晶片（如驍龍8 Gen3、天璣9300）的NPU算力在40-70 TOPS區間，2025-2026年的新一代晶片普遍突破100 TOPS，這是端側跑大模型的硬體門檻。
  \item \textbf{AI Agent場景}：不僅會"問答"，還能跨App執行任務鏈——例如"幫我把昨天開會的照片整理成一個相簿並分享給張三"，系統能自動呼叫相簿、檔案、微信等多個App完成。這是AI手機區別於"帶AI助手的手機"的核心標誌。
\end{itemize}

\booksection{二、各廠商AI戰略對比：誰在認真做，誰在貼牌？}

各家的AI路線差異明顯，選購時可以重點關注：

\begin{itemize}
  \item \textbf{蘋果 Apple Intelligence（iOS 18+）}：強調整合系統級應用——郵件摘要、通知優先順序、寫作潤色、相簿AI修圖。端側模型為主，複雜任務上雲到Private Cloud Compute。優勢是隱私保護和系統流暢度，劣勢是中文場景上線較慢，部分功能2025年才逐步支援。
  \item \textbf{谷歌 Gemini Nano}：Pixel 8/9系列內建，支援"畫圈搜尋"、Magic Compose、錄音摘要等。Gemini Nano是真正意義上的端側小模型，開源生態友好，是安卓AI手機的參考實現。
  \item \textbf{華為盤古大模型}：接入鴻蒙小藝，支援端側文案生成、AI修圖、AI消除（去路人）、通話摘要。鴻蒙NEXT 5.0後進一步系統級整合，AI能力與系統深度融合。
  \item \textbf{小米澎湃大模型（MiLM）}：HyperOS內建，支援AI寫真、AI摘要、AI畫師、超解析度。小米走"端雲協同"路線，端側跑輕量模型，重任務上雲。
  \item \textbf{OPPO（安第斯大模型）與vivo（藍心大模型）}：ColorOS 15與OriginOS 5深度接入，主打AI消除、AI檔案摘要、AI識圖問答、自然語言搜照片（"找去年在海邊拍的照片"）。
\end{itemize}

\textbf{選購建議}：如果你經常需要AI處理工作檔案、整理資訊，優先選端側算力強、AI功能系統級整合的機型；如果只是偶爾用AI修圖，主流中端機已經夠用。

\booksection{三、衛星通訊：從旗艦“噱頭”到中端標配}

衛星通訊是這兩年最實在的技術下沉之一。2022年華為Mate 60和iPhone 14首發時還只是旗艦專屬，2025年已經下放到2000-3000元檔。

\begin{itemize}
  \item \textbf{北斗短報文下放中端機}：北斗衛星的短報文功能（可傳送1000漢字以內的簡訊+位置）已下放到Redmi、realme、iQOO等中端機型，無需月租，免費使用。適合戶外徒步、自駕遊等場景。
  \item \textbf{天通衛星通話}：華為Mate 60/70系列、榮耀Magic6系列等支援天通一號衛星通話功能，可撥打/接聽衛星電話、傳送衛星訊息。需要開通運營商衛星套餐（約每月10-200元），適合經常進無人區的使用者。
  \item \textbf{iPhone衛星訊息（SOS）}：iPhone 14及以後機型支援衛星SOS，可向緊急服務機構傳送文字訊息並共享位置。免費提供2年服務，部分地區續費可用。
  \item \textbf{實用建議}：城市使用者基本用不上，但戶外愛好者、出差到偏遠地區的商務人士強烈建議選帶衛星功能的機型——關鍵時刻能救命。
\end{itemize}

\booksection{四、摺疊屏生態成熟度：能當主力機了嗎？}

摺疊屏已從“科技嚐鮮”走向“主力機可用”。判斷摺疊屏是否成熟，看三點：硬體耐用性、App適配度、互動生態。

\begin{itemize}
  \item \textbf{硬體耐用性}：鉸鏈壽命普遍達50萬次以上（每天開合100次可用13年），水滴鉸鏈減少螢幕摺痕，IPX8級防水已成標配。內屏UTG玻璃+保護膜方案讓螢幕不再"嬌貴"。
  \item \textbf{App適配}：主流App（微信、淘寶、抖音、B站、Office、WPS）已針對大屏做適配，但仍有約30\%的長尾App未適配，顯示為拉伸或黑邊。建議購機前先試用你常用的幾個App。
  \item \textbf{懸停態互動}：摺疊到90°-120°時，部分App支援"上半屏顯示內容、下半屏操作"的懸停模式，適合視訊通話、自拍、看片。相機懸停自拍是摺疊屏獨有的體驗。
  \item \textbf{分屏與多工}：大屏下可同時開兩個App分屏，甚至三個App小窗並存。HyperOS、ColorOS、OriginOS的分屏邏輯已比較成熟，但iOS（iPadOS風格）的分屏在iPhone上仍未開放。
  \item \textbf{價格}：橫折（書本式）從萬元級下探到5000-7000元檔，豎折（翻蓋式）已下探到4000-5000元檔。2025年部分廠商推出3000元檔豎折，摺疊屏真正進入大眾市場。
\end{itemize}

\textbf{選購建議}：如果你是重度閱讀、檔案處理、影片使用者，橫折屏值得考慮；如果你追求便攜和社交屬性，豎折屏更適合。建議至少去線下店把玩30分鐘再決定。

\booksection{五、6G與Wi-Fi 7：下一代連線的速度與邊界}

雖然6G還沒真正落地，但相關技術已經在路上了：

\begin{itemize}
  \item \textbf{6G進展}：3GPP計劃2027年完成6G標準制定，2030年左右商用。6G的理論峰值速率可達1 Tbps（5G的50倍），延遲低至0.1ms，支援"通感一體化"（通訊+感知），手機甚至能當低精度雷達用。目前各廠商還在原型驗證階段，離消費級產品尚遠。
  \item \textbf{Wi-Fi 7（802.11be）}：2024年起成為旗艦機標配，理論峰值速率達46 Gbps（Wi-Fi 6的4倍），支援320MHz頻寬、4K-QAM、MLO（多鏈路同傳）。實測在Wi-Fi 7路由器下，旗艦機傳輸速率可達2-5 Gbps，比Wi-Fi 6E快2-3倍。如果你換了Wi-Fi 7手機，記得配一臺Wi-Fi 7路由器才能發揮全部效能。
  \item \textbf{藍芽5.4與LE Audio}：2024-2026年新機普遍支援藍芽5.4及LE Audio（低功耗音訊），可同時連兩副耳機聽同一首歌，Auracast廣播音訊讓多人在公共場合同步聽同一音源。
\end{itemize}

\booksection{六、矽碳負極電池：續航終於迎來材料級突破}

手機續航多年來被詬病“一天一充”，2024年起矽碳負極電池技術的普及讓情況有了實質改善。

\begin{itemize}
  \item \textbf{技術原理}：傳統石墨負極的理論容量約372 mAh/g，矽的理論容量高達4200 mAh/g（約10倍），但矽充放電時體積膨脹達300\%，容易粉化。矽碳負極透過奈米矽-碳複合結構，把矽的膨脹問題控制在可接受範圍，讓電池能量密度提升15\%-30\%。
  \item \textbf{實際表現}：2024-2025年採用矽碳負極的旗艦機（如小米14 Ultra、榮耀Magic 6至臻版、vivo X100 Pro等），在7-8mm機身內塞進5000-6000 mAh電池，比上代多出約800-1000 mAh，重度使用可撐一天半。
  \item \textbf{百瓦快充配合}：矽碳負極對快充相容性良好，100W/120W快充17分鐘充滿已不是新聞。2025年部分機型支援150W-240W，10分鐘內充至80\%。
  \item \textbf{選購建議}：如果你是重度使用者，優先選採用"矽碳負極+百瓦快充"組合的2024年後機型，續航焦慮會大幅緩解。
\end{itemize}

\booksection{七、本章速覽：一句話記住這些前沿}

\begin{itemize}
  \item AI手機三標準：端側大模型、NPU算力TOPS、AI Agent跨App場景。
  \item 衛星通訊：北斗短報文下放中端，天通通話仍屬旗艦。
  \item 摺疊屏：硬體成熟、App適配七成、懸停/分屏是獨有體驗。
  \item 6G還在路上，Wi-Fi 7已落地，藍芽LE Audio開始普及。
  \item 矽碳負極+百瓦快充：續航與充電的雙突破。
\end{itemize}

\begin{tcbpractice}{來試試：給你的下一臺手機做一份"技術需求清單"}
拿出一張紙，按“必備/加分/不需要”三檔，列出你對AI、衛星、摺疊屏、Wi-Fi 7、快充、電池容量的需求。這樣下次去線下店時，導購推薦的“賣點”你能迅速判斷是真需求還是被推銷。例如：

\begin{itemize}
  \item 必備：5000mAh以上電池、80W以上快充、NFC、IP68防水。
  \item 加分：端側AI摘要、Wi-Fi 7、北斗短報文。
  \item 不需要：摺疊屏（我習慣直板）、衛星通話（不出遠門）、240W極限快充。
\end{itemize}

\end{tcbpractice}

\begin{tcbsummary}{本章小結}
\begin{itemize}
  \item AI手機三大判斷標準：端側大模型、NPU TOPS、AI Agent場景
  \item 各廠商AI戰略：蘋果重隱私整合、谷歌端側開源、華為/小米/OPPO/vivo各有所長
  \item 衛星通訊：北斗短報文下放中端機，天通通話仍屬旗艦專屬
  \item 摺疊屏：硬體已成熟，App適配約70\%，懸停態與分屏是獨有體驗
  \item 6G標準2027年制定，Wi-Fi 7已普及，藍芽LE Audio開始落地
  \item 矽碳負極電池讓5000-6000mAh成為旗艦標配，續航焦慮大幅緩解
\end{itemize}

\end{tcbsummary}
